\section{Marco Teórico}

\subsection{Computación de Alto Rendimiento (HPC)}
La Computación de Alto Rendimiento (HPC, por sus siglas en inglés) consiste en la agregación de potencia de cálculo a través de arquitecturas paralelas para resolver problemas científicos y tecnológicos de extrema complejidad. Esta disciplina permite ejecutar algoritmos y modelados matemáticos que resultarían inabarcables para un sistema tradicional debido a restricciones de tiempo y recursos de hardware \parencite{Hager2010HPC}. Al dividir cargas masivas de trabajo en instrucciones más pequeñas que se procesan simultáneamente, HPC maximiza la velocidad de procesamiento de datos y la obtención de resultados en la ingeniería \parencite{Patterson2017Computer}.

\subsection{Fundamentos Matemáticos de la Multiplicación Matricial}
La multiplicación de matrices es una operación algebraica bidimensional fundamental que procesa dos matrices de origen para producir una tercera. Dadas dos matrices \(A\) de dimensión \(m \times p\) y \(B\) de dimensión \(p \times n\), su producto da como resultado una matriz \(C\) de dimensiones \(m \times n\). Matemáticamente, el valor de cada celda individual \(C_{ij}\) en la matriz resultante se calcula mediante el producto punto o escalar entre la \(i\)-ésima fila de la matriz \(A\) y la \(j\)-ésima columna de la matriz \(B\). Esta operación se expresa de forma analítica mediante la siguiente sumatoria:

\[ C_{ij} = \sum_{k=1}^{p} A_{ik} B_{kj} \]

Para que esta operación esté matemáticamente definida y sea computacionalmente viable, existe una condición estricta de dimensionalidad: el número de columnas de la primera matriz (\(p\)) debe coincidir exactamente con el número de filas de la segunda matriz \parencite{Strang2016Linear}. 

En el contexto de la ingeniería y la Computación de Alto Rendimiento (HPC), la multiplicación matricial trasciende la teoría pura para consolidarse como la rutina principal (\textit{kernel}) subyacente en la gran mayoría del software científico. Esta operación constituye el pilar del Nivel 3 de la especificación BLAS (\textit{Basic Linear Algebra Subprograms}), el estándar de facto para operaciones de álgebra lineal. La eficiencia con la que un procesador ejecuta este algoritmo dicta directamente el rendimiento de aplicaciones críticas modernas, tales como el modelado de dinámica de fluidos, la física de partículas, el procesamiento masivo de imágenes y, muy especialmente, el entrenamiento de arquitecturas profundas en inteligencia artificial mediante redes neuronales \parencite{Dongarra1990BLAS}.


\subsection{Complejidad Computacional de la Multiplicación de Matrices}
La multiplicación de matrices es un problema canónico en el álgebra lineal computacional, utilizado frecuentemente para evaluar el desempeño y la robustez de las arquitecturas de hardware. Dadas dos matrices cuadradas \(A\) y \(B\) de dimensiones \(n \times n\), el algoritmo iterativo secuencial clásico exige iterar mediante tres ciclos anidados, realizando matemáticamente \(n^3\) multiplicaciones y \(n^2(n-1)\) sumas \parencite{Cormen2009Algorithms}. 

Esto establece una complejidad temporal asintótica estricta de \(\mathcal{O}(n^3)\). Este crecimiento cúbico implica que, si el tamaño de la matriz se duplica, el esfuerzo computacional se multiplica por un factor de ocho. Debido a esta pronunciada curva de costo algorítmico, la distribución de la carga mediante técnicas de programación paralela (hilos o procesos) se vuelve un requisito indispensable cuando se operan matrices de dimensiones elevadas \parencite{Grama2003Introduction}.

\subsection{Programación Secuencial y Concurrente}
La programación secuencial es el paradigma clásico de desarrollo donde las instrucciones de un algoritmo se ejecutan de forma estrictamente lineal, una tras otra, sobre un único núcleo de procesamiento. En este modelo tradicional, el límite de rendimiento está directamente dictado por la frecuencia de reloj del hardware. Como respuesta a estas limitaciones físicas, la programación concurrente estructura el software en múltiples hilos o tareas independientes que hacen progreso de forma superpuesta en el tiempo. Cuando estas tareas operan de manera simultánea en una arquitectura de hardware multinúcleo, la concurrencia se traduce en paralelismo físico, lo cual permite dividir la carga de trabajo y reducir drásticamente el tiempo de ejecución en problemas de cómputo intensivo \parencite{BenAri2006Concurrent}.

\subsection{Sistemas de Memoria Compartida}
En el ámbito de la computación paralela, el modelo de memoria compartida establece que múltiples hilos de ejecución (como los creados por la librería Pthreads) operan dentro de un mismo espacio de direccionamiento lógico. A nivel de arquitectura, esto significa que todos los núcleos del procesador pueden acceder de forma directa y simultánea a las mismas estructuras de datos residentes en la memoria principal. Para problemas matriciales, este modelo resulta excepcionalmente eficiente frente a los sistemas de memoria distribuida, ya que evita la necesidad de copiar y transmitir los datos de las matrices de origen repetidamente entre procesos. Aunque este paradigma exige control riguroso para evitar condiciones de carrera, en operaciones intrínsecamente independientes como el cálculo individual de las celdas de una matriz resultado, su implementación maximiza la eficiencia computacional sin sobrecarga de sincronización \parencite{Silberschatz2018Operating}.

\subsection{Jerarquía de Caché, Líneas de Caché y Localidad de Memoria}
El rendimiento de una aplicación de alto rendimiento (HPC) no depende exclusivamente de la capacidad matemática del procesador, sino también del acceso al subsistema de memoria. Para mitigar la alta latencia térmica y física de la memoria RAM, los procesadores modernos integran una jerarquía de memoria ultrarrápida cercana a los núcleos, conocida como memoria Caché, habitualmente dividida en niveles (L1, L2 y L3). La unidad mínima de transferencia de información entre la RAM principal y la memoria caché no es un byte aislado, sino un bloque continuo de datos denominado línea de caché (\textit{cache line}), que en las arquitecturas modernas suele constar de 64 bytes \parencite{Drepper2007Memory}.

Esta transferencia por bloques está fundamentada en el principio de localidad espacial de memoria, un concepto empírico que dicta que si un programa accede a una dirección de memoria, es inminentemente probable que acceda a las direcciones adyacentes a continuación. En lenguajes de programación como C, las matrices bidimensionales se almacenan en la memoria RAM de manera contigua organizadas por filas (\textit{row-major order}). Durante la multiplicación matricial clásica iterativa, la segunda matriz se recorre inevitablemente por columnas; este patrón de acceso ortogonal rompe por completo la localidad espacial. Al intentar leer datos saltando grandes bloques de memoria, los valores necesarios no se encuentran en la línea de caché cargada, provocando fallos de caché (\textit{cache misses}) que obligan al procesador a detenerse (\textit{stall}) mientras busca los datos nuevamente en la lenta memoria RAM, lo cual penaliza severamente el desempeño del algoritmo \parencite{Patterson2017Computer}.

\subsection{Optimización Algorítmica y Aprovechamiento de la Caché}
Aunque la paralelización distribuye la carga de trabajo, el cuello de botella de la memoria solo puede resolverse reestructurando el algoritmo subyacente. La estrategia de optimización más directa para la multiplicación de matrices en C consiste en el intercambio de bucles (\textit{loop interchange}). Altera el orden tradicional de los ciclos anidados de \(i, j, k\) a \(i, k, j\). Matemáticamente, el resultado es idéntico, pero a nivel de hardware, el bucle más interno ahora recorre las matrices \(B\) y \(C\) estrictamente por filas. Esto garantiza que el procesador consuma por completo los 64 bytes de cada línea de caché (\textit{cache line}) cargada antes de solicitar la siguiente, aprovechando al máximo la localidad espacial y reduciendo las penalizaciones por fallos de caché en órdenes de magnitud \parencite{Lam1991Cache}.

Para matrices masivas que exceden la capacidad total de la memoria caché (L1 y L2), la literatura de HPC emplea una técnica avanzada denominada multiplicación por bloques o \textit{Loop Tiling}. Este método particiona las matrices originales en submatrices (o bloques) de un tamaño cuidadosamente calculado para que residan por completo dentro de la caché rápida del procesador. En lugar de multiplicar filas y columnas completas, el algoritmo opera sobre estos bloques aislados, reutilizando los datos cargados repetidas veces antes de que el controlador de memoria los descarte (\textit{eviction}). Esta técnica transforma un problema limitado por el ancho de banda de la memoria (\textit{memory-bound}) en uno puramente computacional (\textit{compute-bound}) \parencite{Grama2003Introduction}.

\subsection{Métricas de Rendimiento en Computación Secuencial}
El análisis de un algoritmo secuencial requiere establecer una línea base absoluta sobre la cual evaluar cualquier mejora algorítmica o de hardware. La métrica fundamental empírica es el Tiempo de Ejecución (\textit{Wall-clock time}), que representa el tiempo físico real transcurrido desde el inicio hasta la finalización del programa. Sin embargo, dado que este valor fluctúa dependiendo de la interferencia del sistema operativo y los accesos a memoria, la arquitectura de computadores recurre a métricas de rendimiento intrínseco (\textit{throughput}) \parencite{Patterson2017Computer}.

Para cargas de trabajo de alto rendimiento algebraico como la multiplicación de matrices, la métrica estandarizada predominante es el volumen de Operaciones de Punto Flotante por Segundo (FLOPS, por sus siglas en inglés). Esta medida cuantifica el rendimiento matemático bruto del procesador y se define mediante la siguiente relación:

\begin{equation}
    \text{FLOPS} = \frac{\text{Operaciones Aritméticas Totales}}{T_{\text{ejecución}}}
\end{equation}

Para el algoritmo clásico iterativo que multiplica dos matrices cuadradas de dimensión \(n \times n\), el número total de operaciones de punto flotante se establece analíticamente en \(2n^3\) (contabilizando una instrucción de multiplicación y una de suma por cada iteración del bucle más interno). Al contrastar los FLOPS experimentales obtenidos en la ejecución secuencial contra el rendimiento teórico máximo (\textit{Peak FLOPS}) que el fabricante del procesador garantiza para un solo núcleo, es posible diagnosticar ineficiencias de memoria o cuellos de botella antes de recurrir a la paralelización \parencite{Hager2010HPC}. 

Adicionalmente, el rendimiento de la ejecución en un único núcleo se evalúa a nivel microarquitectónico mediante el IPC (Instrucciones por Ciclo de reloj). Esta métrica indica la eficiencia con la que el código compilado logra mantener ocupada la segmentación (\textit{pipeline}) del procesador, minimizando las detenciones o burbujas producidas por dependencias de datos en la ejecución lineal \parencite{Yi2020CPI}.


\subsection{Métricas de Rendimiento en Computación Paralela}
Para cuantificar el éxito de una implementación paralela y su eficiencia frente a la ejecución secuencial, la literatura científica emplea métricas matemáticas estandarizadas. La métrica fundamental es el \textit{Speedup} o Aceleración (\(S_p\)), la cual mide la ganancia de velocidad relativa. Se define formalmente como:

\begin{equation}
    S_p = \frac{T_1}{T_p}
\end{equation}

donde \(T_1\) representa el tiempo de ejecución del algoritmo secuencial óptimo sobre un único núcleo, y \(T_p\) es el tiempo de ejecución del algoritmo paralelo empleando \(p\) unidades de procesamiento (hilos o procesos). Derivada de esta métrica se encuentra la Eficiencia (\(E_p\)), que evalúa el grado de aprovechamiento de los recursos físicos, indicando qué fracción del tiempo los procesadores realizan trabajo útil sin quedar ociosos:

\begin{equation}
    E_p = \frac{S_p}{p} = \frac{T_1}{p \cdot T_p}
\end{equation}

El desempeño asintótico de estos sistemas está regido por dos leyes analíticas que abordan la escalabilidad desde diferentes perspectivas teóricas. La \textbf{Ley de Amdahl} evalúa la \textit{escalabilidad fuerte}, demostrando que el \textit{Speedup} máximo de un sistema está estrictamente limitado por la fracción de código que no se puede paralelizar. Asumiendo un tamaño de problema fijo, la ley se expresa como:

\begin{equation}
    S_{\text{Amdahl}} \leq \frac{1}{(1 - f) + \frac{f}{p}}
\end{equation}

donde \(f\) es la fracción del programa que es paralelizable y \((1 - f)\) es la porción estrictamente secuencial (como la inicialización de memoria mediante \texttt{malloc} o la lectura/escritura en disco). Según este postulado matemático, incluso si el número de procesadores \(p\) tiende a infinito, la aceleración máxima nunca superará \(1 / (1 - f)\). Esto explica el rendimiento decreciente observado cuando el uso de múltiples hilos lógicos en un procesador alcanza su límite físico y no aporta mejoras significativas \parencite{Patterson2017Computer}.

En contraste, la \textbf{Ley de Gustafson-Barsis} describe la \textit{escalabilidad débil}. Esta ley argumenta que, en la práctica profesional, los ingenieros no buscan resolver un problema fijo más rápido, sino que utilizan el mayor poder de cómputo para resolver problemas de un tamaño sustancialmente mayor en el mismo marco de tiempo \parencite{Gustafson1988Reevaluating}. Su modelo se define como:

\begin{equation}
    S_{\text{Gustafson}} = (1 - f) + p \cdot f
\end{equation}

Bajo la visión de Gustafson, a medida que la dimensión de las matrices crece de forma masiva, la cantidad de operaciones paralelas \(f\) (multiplicaciones flotantes) aumenta drásticamente, haciendo que la fracción secuencial inicial \((1 - f)\) pierda peso relativo. Esto mitiga el cuello de botella pronosticado por Amdahl, permitiendo alcanzar un \textit{Speedup} escalable y un uso altamente eficiente de los recursos al ajustar la carga de trabajo proporcionalmente al hardware.

\subsection{OpenMP como Modelo de Programación Paralela}
OpenMP (\textit{Open Multi-Processing}) es una interfaz de programación de aplicaciones
(API) de estándar abierto para la programación paralela en sistemas de memoria compartida,
desarrollada y mantenida por el consorcio OpenMP Architecture Review Board. El estándar
define un conjunto de directivas de compilador, rutinas de biblioteca en tiempo de ejecución y
variables de entorno que permiten al programador expresar paralelismo a nivel de bucles y
tareas directamente sobre un código fuente C, C++ o Fortran existente, sin necesidad de
reescribir la arquitectura del programa \parencite{Dagum1998OpenMP}. Esta característica lo
convierte en el modelo predominante para la paralelización incremental en HPC sobre
arquitecturas multinúcleo de memoria compartida.

El modelo de ejecución de OpenMP se fundamenta en el paradigma de bifurcación y unión
(\textit{fork-join}). El programa inicia su ejecución como un hilo único denominado hilo maestro,
el cual ejecuta el código secuencial de forma ordinaria. Al encontrar una directiva de región
paralela, el hilo maestro bifurca un equipo de hilos de trabajo que ejecutan el bloque de código
encerrado de forma concurrente; al concluir dicha región, todos los hilos se sincronizan en una
barrera implícita y se unen de nuevo al hilo maestro, que retoma la ejecución secuencial
\parencite{Dagum1998OpenMP}. Este ciclo puede repetirse tantas veces como regiones paralelas
declare el programa, lo que permite una adopción gradual y controlada del paralelismo sobre
una base de código secuencial ya validada.

La directiva central para la paralelización de bucles es \texttt{\#pragma omp parallel for},
que distribuye automáticamente las iteraciones del bucle entre los hilos del equipo. El
programador puede controlar la política de distribución de trabajo mediante la cláusula
\texttt{schedule}, que admite los modos estático (\textit{static}), dinámico (\textit{dynamic})
y guiado (\textit{guided}), cada uno con implicaciones distintas sobre el balance de carga y la
sobrecarga de coordinación entre hilos \parencite{Chapman2008OpenMP}. Para la
multiplicación de matrices, el modo estático resulta óptimo dado que las iteraciones presentan
una carga de trabajo homogénea; esto permite al compilador asignar bloques contiguos de
filas a cada hilo sin incurrir en la sobrecarga de planificación dinámica en tiempo de ejecución.

El modelo de datos de OpenMP clasifica las variables del programa en dos categorías: compartidas
(\textit{shared}), accesibles y modificables por todos los hilos del equipo simultáneamente, y
privadas (\textit{private}), de las cuales cada hilo posee una copia independiente. El control
explícito sobre este ámbito de visibilidad es un requisito de corrección del programa paralelo:
las matrices de entrada y la matriz resultado deben declararse como compartidas para que todos
los hilos lean y escriban sobre la misma estructura de datos, mientras que las variables de
índice de los bucles internos deben ser privadas para evitar condiciones de carrera
\parencite{Chapman2008OpenMP}. Dado que cada celda \(C_{ij}\) de la matriz resultado es
calculada de forma completamente independiente por un único hilo, la multiplicación de
matrices no requiere mecanismos de sincronización explícita como secciones críticas o
operaciones atómicas, lo que elimina la principal fuente de sobrecarga en implementaciones
OpenMP y permite un escalado cercano al teórico.

\subsection{Perfilado de Código}
El perfilado de código es la disciplina de la ingeniería de software que mide empíricamente el
comportamiento dinámico de un programa en ejecución, con el objetivo de identificar con
precisión cuáles funciones o regiones del código concentran el mayor consumo de tiempo de
CPU o de recursos de memoria. En el contexto de HPC, esta técnica constituye el primer paso
metodológico indispensable antes de cualquier intento de optimización o paralelización, ya que
evita el anti-patrón de optimizar secciones del código que no representan cuellos de botella
reales \parencite{Hager2010HPC}.

Los perfiladores operan bajo dos paradigmas fundamentales. La primera estrategia es la
\textit{instrumentación}, que consiste en insertar sondas de medición en puntos específicos del
código fuente o del binario compilado; este enfoque proporciona datos exactos pero puede
introducir una sobrecarga de ejecución significativa. La segunda estrategia es el
\textit{muestreo estadístico}, en la cual la herramienta interrumpe la ejecución del proceso a
intervalos regulares para registrar el estado del contador de programa; este método produce
resultados estadísticamente representativos con una perturbación mínima sobre el tiempo de
ejecución real \parencite{Grama2003Introduction}.

Entre las herramientas más ampliamente adoptadas en entornos HPC sobre Linux se destacan
dos: gprof (GNU Profiler) y perf. gprof opera combinando instrumentación en tiempo de
compilación con muestreo temporal; analiza el binario para generar un reporte jerarquizado que
muestra el tiempo acumulado por función y el grafo de llamadas, permitiendo identificar con
granularidad de función dónde el programa concentra la mayor fracción de su tiempo de
cómputo \parencite{GrahamEtAl1982Gprof}. perf, por su parte, es una herramienta de análisis
integrada al núcleo Linux que accede directamente a los contadores de hardware del procesador
mediante la interfaz de eventos de rendimiento del sistema operativo. Esto le permite recolectar
métricas microarquitectónicas de bajo nivel, como la tasa de fallos de caché L1/L2/L3, el
número de predicciones erróneas de salto, el IPC real y el número de instrucciones retiradas
por ciclo; datos que resultan imposibles de obtener mediante perfilado por instrumentación de
código fuente \parencite{Patterson2017Computer}.


\subsection{Benchmarks de Rendimiento en HPC}
Un benchmark de rendimiento es un procedimiento de medición estandarizado y reproducible que
ejecuta una carga de trabajo representativa sobre un sistema con el propósito de obtener
métricas comparables entre diferentes arquitecturas de hardware, configuraciones de software
o implementaciones algorítmicas. Para que un benchmark posea validez científica, debe cumplir
criterios de representatividad (que el problema resuelto sea relevante para las aplicaciones de
interés), reproducibilidad (resultados consistentes bajo las mismas condiciones) y equidad (que
no favorezca artificialmente ninguna arquitectura específica) \parencite{Dongarra1990BLAS}.

En el dominio de HPC, el benchmark de referencia histórico y más influyente es LINPACK
(\textit{Linear Algebra Package}), desarrollado por Jack Dongarra, y su sucesor de alta precisión
HPL (\textit{High Performance LINPACK}). El HPL resuelve un sistema denso de ecuaciones
lineales de la forma \(Ax = b\) de dimensión \(N \times N\) mediante la factorización LU con
pivoteo parcial, reportando el rendimiento efectivo en GFLOPS (\(R_{\max}\)) y
contrastándolo contra el rendimiento teórico máximo del hardware (\(R_{\text{peak}}\)). La
eficiencia HPL se define como el cociente \(R_{\max} / R_{\text{peak}}\), y constituye el
criterio de clasificación oficial de la lista TOP500, el ranking mundial de supercomputadoras
\parencite{Dongarra1990BLAS}. Para el benchmarking de operaciones de álgebra lineal densa a
nivel de núcleo computacional, el estándar utilizado en la literatura científica es DGEMM
(\textit{Double-precision General Matrix Multiply}), que forma parte de la especificación BLAS
Nivel 3 y mide directamente el rendimiento de la multiplicación de matrices en doble precisión,
sirviendo como indicador canónico del rendimiento sostenido de punto flotante de un procesador
\parencite{Patterson2017Computer}.


\subsection{Geekbench como Herramienta de Benchmark Multiplataforma}
Geekbench es una suite de benchmarks multiplataforma desarrollada por Primate Labs, diseñada
para evaluar el rendimiento del procesador y la memoria de un sistema mediante cargas de
trabajo sintéticas representativas de aplicaciones reales del mundo moderno. A diferencia de
los benchmarks orientados exclusivamente al cómputo numérico como HPL o DGEMM, Geekbench
adopta un enfoque holístico que cubre categorías como compresión de datos, procesamiento
de imágenes, renderizado de documentos, compilación de código y cargas de trabajo de
aprendizaje automático, con el objetivo de capturar el comportamiento del procesador frente a
una gama amplia de patrones de acceso a memoria y tipos de instrucciones
\parencite{GeekbenchInternals2024}.

La versión actual, Geekbench 6, organiza la evaluación del procesador en dos secciones
principales: núcleo único y múltiples núcleos. Cada sección se subdivide en cargas de trabajo
de aritmética entera y aritmética de punto flotante. La puntuación compuesta final se calcula
como la media aritmética ponderada de las puntuaciones de subsección, asignando un peso del
65\% a las cargas enteras y el 35\% restante a las de punto flotante. Las puntuaciones están
calibradas contra una línea base de 2,500 puntos correspondiente a un procesador Intel Core
i7-12700, de modo que una puntuación doble representa exactamente el doble del rendimiento
de referencia \parencite{GeekbenchInternals2024}. Geekbench 6 introduce además un modelo de
paralelismo basado en tareas compartidas en sustitución del esquema de tareas independientes
de versiones anteriores, lo cual modela con mayor fidelidad el comportamiento real de las
aplicaciones modernas que distribuyen trabajo entre núcleos \parencite{GeekbenchInternals2024}.

Desde una perspectiva microarquitectónica, el documento técnico oficial de Geekbench 6
caracteriza cada carga de trabajo mediante métricas de bajo nivel que incluyen IPC, tasa de
fallos de predicción de salto, tamaño del conjunto de trabajo activo y tasas de fallos de caché
para los niveles L1D, L2 y L3. Por ejemplo, cargas de trabajo con patrones de acceso
altamente irregulares como la carga de navegación basada en el algoritmo de Dijkstra registran
tasas de fallo en L3 del 23.9\% en modo de núcleo único, mientras que cargas con alta
localidad espacial como la compresión de datos exhiben tasas de fallo L3 del 12.8\%,
evidenciando cómo el comportamiento de caché impacta directamente la puntuación final del
benchmark \parencite{GeekbenchInternals2024}.